\section{Summing all residue families with the exact loss}\label{sec:residue-summation}

\status{The actual degree padding and all numerical and summation inputs to
\texttt{RestSum} are proved here. The segment functions are supplied by GP;
the injection into the coarse residue space is supplied by LA.}

Write $k=r+K$, $N=|TT|$, $B(n)=\binom{n+k-1}{k-1}$, and
$f_{j,i}=\mathrm{fr}(j,i)$. Set $z_j=f_{j,r-1}$ and $s_j=f_{j-1,0}$.
The late letters of segment $j$ are the selected lifts in $(z_{j-1},z_j]$;
they split into $LL_j$ and the current frame. Over one period they tile
$TT$ bijectively. A full segment lies in $(z_{j-2},z_j]$, so a residue
occurs at most twice in the full segment sets. In particular $J\le N$.
These statements follow directly from ordered periodic frames. \textsf{[RS1]}

\paragraph{The actual numerical margins.}
Each determinant row contributes at most $2a/(k^2Y)$ to its residue's
degree. There are at most $N_s$ determinants per segment; the preceding
two-occurrence bound gives
\[
 E_t\le\frac{4N_sa}{k^2Y},\qquad
 a'_t=\deg_t-E'_t-E_t\ge\frac9{10}\deg_t.
 \tag{RS.1}
\]
For the second inequality, the source hypotheses give
$a\ge32$, $k\ge5$, $N_s\ge1$, $k\ge32N_s$, $E_0\le a/32$,
and hence $E_t\le a/40$. They also give the sharper total loss
$E'_t+E_t\le5\deg_t/(kY)$. Thus all these natural subtractions are
ordinary nonnegative differences. \textsf{[RS2]}

Put $U=8Yk^2N_s$, $T_j=a/(Up^{z_j-s_j})$,
$R_j=\log_p\lfloor T_j\rfloor$, and $P_j(x)=p^{z_j-x+R_j}$.
The size assumption gives
\[
 T_j\ge4p(E_0+m_*+1),\quad
 p^{R_j}\le T_j<p^{R_j+1},\quad
 \frac{a'_{\operatorname{pos}x}}{P_j(x)}
 \ge\frac{36}{5}Yk^2N_s. \tag{RS.2}
\]
The scale relation has a short exact proof. With
$H=\log_p\lfloor a/U\rfloor$,
\[
 A_j=z_j+R_j=\max(z_j,H+s_j),\qquad
 P_j(x)=p^{A_j-x},\qquad Q_j=p^{A_{j+1}-A_j}.
\]
Both sequences inside the maximum strictly increase. Since
$A_{j+J}=A_j+n'$, this proves
\[
 P_j(x)Q_j=P_{j+1}(x),\qquad \prod_{j=0}^{J-1}Q_j=p^{n'}=q.
 \tag{RS.3}
\]
It also proves periodicity and consistency at equal-residue late
occurrences. \textsf{[RS3; LA2]}

\paragraph{Degree padding and digit partitions.}
For a label $\lambda_t$ let $L_t=\sum_c\lambda_{t,c}$, with
$0\le\lambda_{t,c}<P$, $L_t\le a'_t$ and $L_t\equiv a'_t\pmod P$.
Then $n_j(x)=(a'_t-L_t)/P_j(x)$ is an integer in
$(a'_t/P_j(x)-k,a'_t/P_j(x)]$. Terminal degrees satisfy
$0\le e_{j,i}<rQ_j$ and $e_{j,i}\equiv n_j(f_{j,i})\pmod{Q_j}$.
Equation (RS.2) in segment $j+1$ gives $n_j(f_{j,i})>e_{j,i}$.
The old-frame degree is therefore exactly $(n_j-e_{j,i})/Q_j$.

More generally, if $s<kP$ and $e<rQ$, successive floors give
\[
 \left\lfloor\frac{(\lfloor(n-s)_+/P\rfloor-e)_+}{Q}\right\rfloor
 \ge\left\lfloor\frac n{PQ}\right\rfloor-k-r.
\]
This proves the source's weaker old-frame margin $k+r+1$ and hence
\texttt{RestPad}. The exact decomposition $w_c=\lambda_c+Pv_c$ gives
\[
 B(n)=\sum_{\substack{0\le\lambda_c<P,\ L\le n\\L\equiv n\pmod P}}
 B\!\left(\frac{n-L}{P}\right).
 \tag{RS.4}
\]
If $kP\le n$, there are exactly $P^{k-1}$ labels satisfying that
congruence: any $k-1$ coordinates determine the last. \textsf{[RS4--RS5]}

\paragraph{Two binomial comparisons.}
The product formula for $B$ and Bernoulli's inequality show
\[
\begin{split}
 B(a'_t)&\ge(1-5/Y)B(\deg_t),\\
 (P')^{k-1}B(b)&\ge(1-1/Y)B(a'_t),
 \quad b=\bigl(\lfloor a'_t/P'\rfloor-(k+r+1)-(kN_s+1)\bigr)_+,
\end{split}\tag{RS.5}
\]
where the second line uses $P'=P_j(f_{j,i})Q_j$.
Indeed (RS.1) loses at most $5\deg_t/(kY)$, while
$P'b\ge a'_t-4kN_sP'\ge(1-1/(kY))a'_t$ by (RS.2).
Compare each of the $k-1$ factors and use $(1-\theta)^{k-1}\ge1-(k-1)\theta$.
\textsf{[RS6]}

\paragraph{The exact summation.}
Let $\alpha=1-a^{-N_s^2}$ and $\tau=\max(0,1-r/(2Y))$.
Suppose $G_j(b)\ge0$ depends only on late degrees and obeys the source's
terminal-residue lower bound
\[
 \sum_{e_j}G_j(\operatorname{setTau}_j(b,e_j))
 \ge\alpha W_j[Q_j^{r-1}\tau]^r\prod_{x\in LL_j}B(b(x)),
\]
where $W_j=\prod_i B((\mathrm{Bw}_j(f_{j-1,i})-kN_s-1)_+)$.
For a fixed label, the choices $e_j$ are independent; finite
distributivity makes the sum of products a product of sums.
Now sum the labels using (RS.4). At an $LL$ residue this supplies
$B(a'_t)$; at a frame residue it supplies $P_j(f_{j,i})^{k-1}$.
Periodicity shifts $W_j$ to $W_{j+1}$. Multiplication by $q^{rK}$
supplies $Q_j^K$ at each of the $r$ frame entries, and
\[
 Q_j^{r-1}P_j(f_{j,i})^{k-1}Q_j^K
   =(P_j(f_{j,i})Q_j)^{k-1}.
\]
Thus (RS.5), the tiling, and $\eta=(1-5/Y)(1-1/Y)(1-r/(2Y))$ give
\[
 q^{rK}\sum_{\lambda\in LA}\sum_{e\in EE(\lambda)}
       \prod_jG_j(b_j(\lambda,e))
 \ge\alpha^J\eta^N\prod_{t\in TT}B(\deg_t).
 \tag{RS.6}
\]
If $1-Nr^2k/Y>0$, all factors used above are nonnegative,
$\alpha\ge1-1/(32Y)$, and SD7 gives
$\alpha^J\eta^N\ge1-Nr^2k/Y$.
A nonpositive target follows directly from the nonnegative sum.
This proves \texttt{RestSum 1088 1} with every actual summation and
numerical premise discharged. \textsf{[RS7--RS8]}
